\SGASubsection{4.3}{Opposite parabolic subgroups}
\SGARefTarget{sga3:xxvi:subsection:4:3}

\SGAInlineRefTarget{sga3:xxvi:paragraph:4-3-1}\paragraph{\SGARefLink{sga3:xxvi:paragraph:4-3-1}{4.3.1}.}
Let \(G\) be an \(S\)-reductive group.  Exp.\ XXIV, \SGARefLink{sga3:xxiv:proposition:3-16-6}{3.16.6} defines a
canonical ``outer automorphism'' \(s_G\) of \(G\), of order at most
\(2\).\footnote{\emph{Editor's note.} ``Of order \(2\)'' was replaced
by ``of order at most \(2\)'', since \(s_G\) may be trivial, for example
when \(G\) has type \(A_1,B_n,C_n,\ldots\).}  It induces a canonical
automorphism, again of order at most \(2\), of
\(\underline{\operatorname{Dyn}}(G)\), and hence of
\(\Of(\underline{\operatorname{Dyn}}(G))\).  We denote all of these
automorphisms by \(s_G\), or simply by \(s\).  Two types of parabolic
subgroups
\[
 t,t'\in\Of(\underline{\operatorname{Dyn}}(G))(S)
\]
are said to be \emph{opposite} if \(t=s_G(t')\).
\SGARefTarget{sga3:xxvi:item:4:3:1}

\begin{sgastatement}{Theorem 4.3.2}
Let \(S\) be a scheme, \(G\) an \(S\)-reductive group, and \(P\) a
parabolic subgroup of \(G\).

\textup{(a)} If \(L\) is a Levi subgroup of \(P\), there is a unique
parabolic subgroup \(P'\) of \(G\) such that \(P\cap P'=L\).

\textup{(b)} For every parabolic subgroup \(Q\) of \(G\), the following
conditions are equivalent:
\begin{enumerate}[leftmargin=3.2em,labelsep=.7em]
\item[\textup{\SGAInlineRefTarget{sga3:xxvi:item:4:3:1:item:i}(i)}]
For every \(s\in S\), the identity component
\(\bigl((P\cap Q)_{\bar s}\bigr)^0\), which is smooth by \SGARefLink{sga3:xxvi:lemma:4:1:1}{4.1.1}, is
reductive.
\item[\textup{\SGAInlineRefTarget{sga3:xxvi:item:4:3:1:item:ii}(ii)}]
The group \(P\cap Q\) is a Levi subgroup of both \(P\) and \(Q\).
\item[\textup{\SGAInlineRefTarget{sga3:xxvi:item:4:3:1:item:iii}(iii)}]
The groups \(P,Q\) have opposite types and the pair \((P,Q)\) is in
transversal position.
\item[\textup{\SGAInlineRefTarget{sga3:xxvi:item:4:3:1:item:iv}(iv)}]
The groups \(P,Q\) have opposite types and
\(\rad^u(P)\cap Q=e\).
\item[\textup{\SGAInlineRefTarget{sga3:xxvi:item:4:3:1:item:v}(v)}]
\[
 \rad^u(P)\cap Q=\rad^u(Q)\cap P=e.
\]
\item[\textup{\SGAInlineRefTarget{sga3:xxvi:item:4:3:1:item:vi}(vi)}]
The canonical morphism \(\rad^u(P)\times_SQ\to G\) is an open immersion.
\item[\textup{(vi\(^{\prime}\))}]
The canonical morphism \(\rad^u(P)\to G/Q\) is an open immersion.
\item[\textup{\SGAInlineRefTarget{sga3:xxvi:item:4:3:1:item:vii}(vii)}]
There is an étale covering family \(\{S_i\to S\}\) and, for every \(i\),
a splitting \((T_i,M_i,R_i)\) of \(G_{S_i}\) and a subset
\(R_i^{(1)}\subset R_i\), such that \(P_{S_i}\) is the subgroup of type
\textup{(R)} containing \(T_i\) and defined by \(R_i^{(1)}\), while
\(Q_{S_i}\) is the subgroup defined by
\[
 R_i^{(2)}=-R_i^{(1)}.
\]
\end{enumerate}
\end{sgastatement}
\SGARefTarget{sga3:xxvi:theorem:4:3:2}

\SGAProof
We first prove part \textup{(b)}.  Conditions \textup{\SGARefLink{sga3:xxvi:theorem:4-2-1:item:iii}{(iii)}} and
\textup{\SGARefLink{sga3:xxvi:theorem:4-2-1:item:vii}{(vii)}} are equivalent by \SGARefLink{sga3:xxvi:theorem:4-2-1}{4.2.1}\textup{\SGARefLink{sga3:xxvi:theorem:4-2-1:item:vii}{(vii)}} and the
definition of \(s_G\) in the split case (Exp.\ XXIV,
\SGARefLink{sga3:xxiv:proposition:3-16-2}{3.16.2}\textup{\SGARefLink{sga3:xxiv:proposition:3-16-2:item:iv}{(iv)}}).  Plainly \textup{\SGARefLink{sga3:xxiv:proposition:3-16-2:item:ii}{(ii)}} implies \textup{\SGARefLink{sga3:xxiv:proposition:3-16-2:item:i}{(i)}}, and
\textup{(vi\(^{\prime}\))} implies \textup{\SGARefLink{sga3:xxvi:theorem:4-2-1:item:vi}{(vi)}} by base change along
\(G\to G/Q\).

Suppose now that \textup{\SGARefLink{sga3:xxvi:item:4:3:1:item:vii}{(vii)}} holds.  Since every assertion to be
proved is étale-local, assume that \(G=(G,T,M,R)\) is split, that \(P\)
is defined by \(R'\subset R\), and that \(Q\) is defined by \(-R'\).
Let \(L\) be the subgroup of type \textup{(R)} containing \(T\) and
defined by \(R'\cap(-R')\).  By Exp.\ XXII, \SGARefLink{sga3:xxii:lem:5.11.3}{5.11.3}, \(L\) is a Levi
subgroup of both \(P\) and \(Q\).  Moreover,
\[
 P=L\cdot\rad^u(P),\qquad Q=L\cdot\rad^u(Q),
\]
and Exp.\ XXII, \SGARefLink{sga3:xxii:cor:5.6.7}{5.6.7} gives
\[
 \rad^u(P)\cap Q=Q\cap\rad^u(P)=e.
\]
Thus \(P\cap Q=L\), proving \textup{\SGARefLink{sga3:xxii:cor:5.6.7:item:ii}{(ii)}} and \textup{\SGARefLink{sga3:xxvi:item:4:3:1:item:v}{(v)}}.

Since \(P,Q\) are in transversal position, the canonical morphism
\(P\to G/Q\) induces an open immersion
\[
 P/(P\cap Q)\longrightarrow G/Q
\]
by \SGARefLink{sga3:xxvi:theorem:4-2-1}{4.2.1}.  The canonical morphism
\[
 \rad^u(P)\longrightarrow P/(P\cap Q)=P/L
\]
is an isomorphism, so \textup{(vi\(^{\prime}\))} follows.  Hence all
the conditions follow from \textup{\SGARefLink{sga3:xxvi:theorem:4-2-1:item:vii}{(vii)}}.

It remains to show that each of \textup{\SGARefLink{sga3:xxvi:item:4:3:1:item:i}{(i)}}, \textup{\SGARefLink{sga3:xxvi:item:4:3:1:item:iv}{(iv)}},
\textup{\SGARefLink{sga3:xxvi:item:4:3:1:item:v}{(v)}}, and \textup{\SGARefLink{sga3:xxvi:item:4:3:1:item:vi}{(vi)}} implies \textup{\SGARefLink{sga3:xxvi:item:4:3:1:item:vii}{(vii)}}.  Since the
equivalence of \textup{\SGARefLink{sga3:xxvi:item:4:3:1:item:ii}{(ii)}} and \textup{\SGARefLink{sga3:xxvi:item:4:3:1:item:iii}{(iii)}} is already known, this
may be checked on geometric fibers.  We may suppose that \(S\) is the
spectrum of an algebraically closed field.  By \SGARefLink{sga3:xxvi:lemma:4:1:1}{4.1.1}, a maximal torus
\(T\) of \(G\) is contained in \(P\cap Q\).  Let \(R,R_1,R_2\) be the
sets of roots of \(G,P,Q\), respectively, relative to \(T\).  Write
\[
 R_1^a=\{\alpha\in R_1\mid-\alpha\notin R_1\},
\]
and define \(R_2^a\) similarly.  We must prove \(R_1=-R_2\).

\begin{itemize}[leftmargin=2em]
\item Condition \textup{\SGARefLink{sga3:xxvi:item:4:3:1:item:i}{(i)}} says that \(R_1\cap R_2\) is symmetric.
If \(\alpha\in R_1\) and \(\alpha\notin R_2\), then
\(\alpha\in-R_2\); if \(\alpha\in R_2\), then
\[
 \alpha\in R_1\cap R_2=-(R_1\cap R_2)\subset-R_2.
\]
Thus \(R_1\subset-R_2\), and symmetry gives \(R_1=-R_2\).

\item Condition \textup{\SGARefLink{sga3:xxvi:item:4:3:1:item:iv}{(iv)}} gives
\(\operatorname{Card}(R_1)=\operatorname{Card}(R_2)\) and
\(R_1^a\cap R_2=\varnothing\).  The latter condition is equivalent to
\(R_2\subset-R_1\), and equality of cardinalities then gives
\(R_2=-R_1\).

\item Condition \textup{\SGARefLink{sga3:xxvi:item:4:3:1:item:v}{(v)}} gives
\[
 R_1^a\cap R_2=R_2^a\cap R_1=\varnothing.
\]
Hence \(R_2\subset-R_1\) and \(R_1\subset-R_2\), so again
\(R_2=-R_1\).

\item Condition \textup{\SGARefLink{sga3:xxvi:item:4:3:1:item:vi}{(vi)}} gives
\[
 \operatorname{Lie}(\rad^u(P))\oplus\operatorname{Lie}(Q)
 =\operatorname{Lie}(G).
\]
Thus \(R\) is the disjoint union of \(R_2\) and \(R_1^a\), which implies
\(R_2=-R_1\).
\end{itemize}

This proves part \textup{(b)}.  For part \textup{(a)}, the implication
\textup{\SGARefLink{sga3:xxvi:item:4:3:1:item:vii}{(vii)}} \(\Rightarrow\) \textup{\SGARefLink{sga3:xxvi:item:4:3:1:item:ii}{(ii)}} already gives the desired
parabolic subgroup \(P'\) étale-locally.  It remains only to prove
uniqueness, and this too is étale-local.  We may suppose that \(G\) is
split relative to a maximal torus \(T\) of \(L\), with \(P\) and \(P'\)
defined by subsets \(R_1\) and \(R'_1\) of \(R\).  By hypothesis
\(R_1\cap R'_1\) is symmetric.  The preceding argument gives
\(R'_1=-R_1\); hence \(P'\) is determined by \(P\) and \(L\).
\SGAQED

\begin{sgastatement}{\SGAInlineRefTarget{sga3:xxvi:definition:4-3-3}Definition 4.3.3}
Two parabolic subgroups of \(G\) satisfying the equivalent conditions
\textup{(i)}--\textup{\SGARefLink{sga3:xxvi:item:4:3:1:item:vii}{(vii)}} of \SGARefLink{sga3:xxvi:item:4:3:1}{Theorem 4.3.2} are said to be
\emph{opposite}.  If \(P\) is a parabolic subgroup of \(G\) and \(L\)
is a Levi subgroup of \(P\), the \emph{parabolic subgroup opposite to
\(P\) relative to \(L\)} is the unique parabolic subgroup \(Q\) of
\(G\) such that \(P\cap Q=L\).

If \(T\) is a maximal torus of \(P\), the \emph{parabolic subgroup
opposite to \(P\) relative to \(T\)} is the unique \(Q\) such that
\(P\cap Q\) is the unique Levi subgroup of \(P\) containing \(T\)
\textup{(\SGARefLink{sga3:xxvi:proposition:1-6}{1.6})}; equivalently, it is the unique \(Q\) such that
\(P\cap Q\) contains \(T\) and \(P,Q\) are opposite.
\end{sgastatement}
\SGARefTarget{sga3:xxvi:def:4:3:3}

By \SGARefLink{sga3:xxvi:item:4:3:1}{4.3.2}\textup{\SGARefLink{sga3:xxvi:item:4:3:1:item:iii}{(iii)}}, \SGARefLink{sga3:xxvi:corollary:4-2-4}{Corollaries 4.2.4} and \SGARefLink{sga3:xxvi:corollary:4-2-5}{4.2.5} immediately yield
parallel results.  We record one sample.

\begin{sgastatement}{\SGAInlineRefTarget{sga3:xxvi:corollary:4-3-4}Corollary 4.3.4}
Let \(S\) be a scheme, \(G\) an \(S\)-reductive group, and \(P,Q\) two
parabolic subgroups of \(G\).
\begin{enumerate}[label=\textup{(\roman*)}]
\item The groups \(P,Q\) are opposite if and only if
\(P_{\bar s},Q_{\bar s}\) are opposite for every \(s\in S\).  If
\(S'\to S\) is surjective and \(P_{S'},Q_{S'}\) are opposite, then
\(P,Q\) are opposite.
\item Let \(\Opp(G)(S')\) be the set of pairs of opposite parabolic
subgroups of \(G_{S'}\), for \(S'\to S\).  This functor is represented
by an open subscheme of \(\Par(G)^2\).  Let \(\Opp(/P)(S')\) be the set
of parabolic subgroups of \(G_{S'}\) opposite to \(P_{S'}\).  This
functor is represented by an open subscheme, universally schematically
dense over \(S\), of
\[
 \Par_{s(t(P))}(G).
\]
\item Suppose \(P,Q\) are opposite, and let \(P',Q'\) be parabolic
subgroups of \(G\), with \(P'\) of the same type as \(P\).  The groups
\(P',Q'\) are opposite if and only if, étale-locally, the pair
\((P',Q')\) is conjugate to \((P,Q)\).
\textup{(In §5 the étale topology will be replaced by the Zariski
topology.)}
\end{enumerate}
\end{sgastatement}
\SGARefTarget{sga3:xxvi:cor:4:3:4}
\footnote{\emph{Editor's note.} Here again, ``relatively dense'' was
replaced by ``universally schematically dense over \(S\)''; compare
EGA IV\textsubscript{3}, Definition 11.10.8.}

\begin{sgastatement}{\SGAInlineRefTarget{sga3:xxvi:corollary:4-3-5}Corollary 4.3.5}
Let \(S\) be a scheme, \(G\) an \(S\)-reductive group, and \(P\) a
parabolic subgroup of \(G\).
\begin{enumerate}[label=\textup{(\roman*)}]
\item The morphism
\[
 \Opp(/P)\longrightarrow\Lev(P),\qquad Q\longmapsto P\cap Q,
\]
is an isomorphism.  The scheme \(\Opp(/P)\) is a principal homogeneous
bundle under \(\rad^u(P)\), acting by inner automorphisms.  If \(S\) is
affine, \(G\) has a parabolic subgroup opposite to \(P\).
\item Suppose \(S\) is semilocal, with closed points \(\{s_i\}\).  For
each \(i\), let \(Q_i\) be a parabolic subgroup of \(G_{s_i}\) opposite
to \(P_{s_i}\).  There is a parabolic subgroup \(Q\) of \(G\), opposite
to \(P\), such that \(Q_{s_i}=Q_i\) for every \(i\).
\item The morphism
\[
 \Opp(G)\longrightarrow\mathrm{PL},\qquad
 (P,Q)\longmapsto(P,P\cap Q),
\]
is an isomorphism.
\end{enumerate}
\end{sgastatement}
\SGARefTarget{sga3:xxvi:cor:4:3:5}

These assertions follow from \SGARefLink{sga3:xxvi:item:4:3:1}{Theorem 4.3.2}\textup{(a)} and from \SGARefLink{sga3:xxvi:cor:1:9}{1.9},
\SGARefLink{sga3:xxvi:cor:2:3}{2.3}, and \SGARefLink{sga3:xxvi:cor:2:8}{2.8}.

\begin{sgastatement}{Remark 4.3.6}
Let \(P,Q\) be opposite parabolic subgroups of \(G\), and let
\(P\cdot Q\) be the open subscheme of \(G\), introduced in \SGARefLink{sga3:xxvi:remark:4:2:6}{4.2.6}, which
is the sheaf-theoretic image of \(P\times_SQ\).  The product morphism
\(G\times_SG\to G\) induces isomorphisms
\[
 \rad^u(P)\times_SQ
 \xrightarrow{\sim}P\cdot Q
 \xleftarrow{\sim}P\times_S\rad^u(Q).
\]
\end{sgastatement}
\SGARefTarget{sga3:xxvi:remark:4:3:6}

Indeed, this follows from \SGARefLink{sga3:xxvi:item:4:3:1}{4.3.2} (or \SGARefLink{sga3:xxvi:corollary:4-2-5}{4.2.5}\textup{\SGARefLink{sga3:xxvi:item:4:3:1:item:ii}{(ii)}}) and from the fact
that \(P\cap Q\) is a Levi subgroup of both \(P\) and \(Q\):
\[
 P=\rad^u(P)\cdot(P\cap Q),\qquad
 Q=\rad^u(Q)\cdot(P\cap Q).
\]
Likewise, there is a commutative diagram
\[
\begin{tikzcd}[column sep=large,row sep=large]
 \rad^u(P) \arrow[r,hook] \arrow[d,"\sim"] & G/Q\arrow[d,"\sim"]\\
 \Opp(/P) \arrow[r,hook] & \Par_{t(Q)}(G),
\end{tikzcd}
\]
in which the vertical arrows are induced by
\(g\mapsto\Int(g)Q\).
